\documentclass{article}
\usepackage[margin=0.8in]{geometry}
\usepackage{amssymb}

\title{Lectures on the Geometric Anatomy of\\Theoretical Physics}
\author{Frederic P. Schuller}
\date{}

\begin{document}
\maketitle
\newpage
\tableofcontents

\section*{Introduction}
Theoretical physics is about casting our concept of nature into rigorous mathematical form for better and worse. Theoretical physics, however, does not do this for its own sake. It does so in order to fully explore the implications of what our concept about the real world are. So to certain extend, the spirit of theoretical physics can be cast into the words of Wittgenstein who said, ``What we cannot speak about clearly, that we must pass over in silence.'' If we have concept about real world and if it is not possible to cast them in precise mathematical form, than it is usually an indicator that some aspect of these concepts have not been well understood. Theoretical physics aims at casting these concepts to mathematical language. But again mathematics is just that; a language. If we want to extract physical conclusions from this formulation, we must interpret the language. This interpretation is not the task of mathematics but that of physics. To again use Wittgenstein, he says, ``The theorems of mathematics all say the same, namely \textit{nothing}.'' Obviously, he does not mean that mathematics is useless, he refers to the fact that if you have a theorem of type $A$ if and oly if $B$ ($A$ and $B$ being propositions), then, obviously, $B$ says nothing else than $A$ does vice versa. Mathematically, or logically, it is a tautology but psychologically for our understanding of $A$, it may be very useful to have reformulation of $A$ in terms of $B$. With this view, that mathematics is language, the idea of this course is to provide proper language for theoretical physics. 
 
\subsection*{Aim of the course}
\subsection*{Structure of the course}

\section{Logic of propositions and predicates}
\section{Axioms of set theory}
\section{Classification of Sets}
\section{Topological spaces: Construction and purpose}
\section{Topological spaces: Some heavily used invariants}
\section{Topological manifolds and manifold Bundles}
\section{Differentiable structures: Definition and classification}
\section{Tensor Theory I: Over a Field}
\section{DIfferential structures: The pivotal concept of tangent vector spaces}
\section{Construction of the tangent bundle}
\section{Tensor Space Theory II: Over a Ring}
\section{Grassmann Algebra and deRahm Cohomology}
\section{Lie Groups and their Lie Algebras}
\section{Classification of Lie algebras and Dynkin diagrams}
\section{The Lie Group $SL(2,\mathbb{C})$ and Its Lie Algebra $SL(2,\mathbb{C})$}
\section{Dynkin diagrams from Lie algebras, and vice versa}
\section{Representation theory of Lie groups and Lie algebras}
\section{Reconstruction of a Lie group from its algebra}
\section{Principle fibre bundles}
\section{Associated fibre bundles}
\section{Connections and connection 1-forms}
\section{Local representations of a connection on the base manifold: Yang-Mills fields}
\section{Parallel transport}
\section{Curvature and torsion on principle bundles}
\section{COvariant derivatives}
\section{Application: Quantum mechanics on curved spaces}
\section{Application: Spin structures}
\section{Application: Kinematical and dynamical symmetries}

\end{document}
